\section*{الفصل \ref{ch:g9:microbes-and-infection} --- الميكروبات والعدوى}

\begin{solution}{exo:g9:microbes-and-infection:1}
\omterm{def:g9:microbes-and-infection:sorted}{البكتيريا} (\omterm{def:g5:first-look-at-cells:cell}{خلايا} مفردة بلا \omterm{def:g6:cell-unit-of-life:parts}{نواة})، \omterm{def:g6:decomposers-and-soil:decomposers}{والفطريات} وحيدة \omterm{def:g5:first-look-at-cells:cell}{الخلية}
\omterm{ex:g6:producing-our-food:bread}{كالخميرة} وسائر حياة البركة وحيدة \omterm{def:g5:first-look-at-cells:cell}{الخلية} --- \omterm{def:g9:microbes-and-infection:sorted}{والفيروسات}،
وهي مختلفة في \omterm{def:g5:biodiversity:species}{النوع}: لا \omterm{def:g5:first-look-at-cells:cell}{خلية} لها ولا حياة خاصة بها، بل
نص وراثي معبأ يستعير آلة \omterm{def:g5:first-look-at-cells:cell}{خلية} حية.
\end{solution}

\begin{solution}{exo:g9:microbes-and-infection:2}
تريليونات \omterm{def:g9:microbes-and-infection:sorted}{البكتيريا} غير الضارة المقيمة على \omterm{def:g1:the-five-senses:sense}{الجلد} وفي
المعي. والخدمتان: إعانة \omterm{def:g4:journey-of-food:digestion}{الهضم}، واحتلال الأرض --- فتزاحم
\omterm{prop:g9:microbes-and-infection:mostly}{الممرضات} عن مواضع الهبوط.
\end{solution}

\begin{solution}{exo:g9:microbes-and-infection:3}
الانتقال (رذاذ سعلة)؛ والدخول (جرح، أو المسالك
الهوائية)؛ والتضاعف (انقسامات تضاعف العدد كل ساعة، أو
\omterm{def:g5:first-look-at-cells:cell}{خلايا} مخطوفة تنسخ)؛ والمرض (خفقان الإصبع المتقيح، وخشونة
حلق الزكام).
\end{solution}

\begin{solution}{exo:g9:microbes-and-infection:4}
\omterm{def:g9:microbes-and-infection:sorted}{بكتيريا} الجرح تتضاعف بين \omterm{def:g5:first-look-at-cells:cell}{الخلايا} في النسيج الدافئ الرطب،
فتسمم جوارها. \omterm{def:g9:microbes-and-infection:sorted}{وفيروس} الزكام يتضاعف داخل \omterm{def:g5:first-look-at-cells:cell}{خلايا} بطانة
المسالك الهوائية، وكل \omterm{def:g5:first-look-at-cells:cell}{خلية} مخطوفة تطلق حشدًا من النسخ.
\end{solution}

\begin{solution}{exo:g9:microbes-and-infection:5}
\omterm{def:g1:the-five-senses:sense}{الجلد}: السور عند الدخول. ومكانس البطانات وحمّاماتها
(تنظيف المسالك الهوائية، والدمع، وحمض \omterm{def:g4:journey-of-food:tract}{المعدة}): الدخول
والتضاعف الباكر. \omterm{prop:g9:microbes-and-infection:mostly}{والنبيت المقيم}: الدخول --- فالأرض
محتلة. وصابون \omterm{def:g1:healthy-habits:hygiene}{النظافة} والطهو والتبريد: الانتقال
والتضاعف.
\end{solution}

\begin{solution}{exo:g9:microbes-and-infection:6}
الانتقال --- مسافر الحمى محمولًا على أيدي الأطباء من
غرفة التشريح إلى جناح الوضع. وباستور بيّن المسافرين
أنفسهم: \omterm{def:g6:producing-our-food:fermentation}{الميكروبات} عوامل التحمّض والمرض \omterm{def:g9:microbes-and-infection:stages}{والعدوى} ---
والحرارة موقفًا لها.
\end{solution}

\begin{solution}{exo:g9:microbes-and-infection:7}
المضادات الحيوية تجوّع \omterm{def:g9:microbes-and-infection:sorted}{البكتيريا} وتكسرها --- فهي \omterm{def:g5:first-look-at-cells:cell}{خلايا}
مستقلة لها آلتها التي تسمَّم. أما \omterm{def:g9:microbes-and-infection:sorted}{فيروس} الزكام فلا آلة
له: إنه يجري على آلتك \emph{أنت}، والدواء الذي يصوّب إلى
عمله يصوّب إلى عمل خلاياك --- ولهذا يترك دواء \omterm{def:g9:microbes-and-infection:sorted}{البكتيريا}
\omterm{def:g9:microbes-and-infection:sorted}{الفيروس} من غير مساس.
\end{solution}

\begin{solution}{exo:g9:microbes-and-infection:8}
اغسل قبل الأكل: يسد الانتقال من اليد إلى الفم. واغسل بعد
المرحاض: يسد حلقة الخروج والدخول في طريق المعي. وغطِّ
السعلة (وعطسة \omterm{def:g1:my-body:joint}{المرفق}): تسد انتقال الرذاذ عند المنبع.
(وقواعد الأسنان والطعام تسد مواضع التضاعف في الفم.)
\end{solution}

\begin{solution}{exo:g9:microbes-and-infection:9}
الممرض: \omterm{def:g9:microbes-and-infection:sorted}{بكتيريا} التسمم الغذائي، على الدجاج. والانتقال:
الطبق نفسه. والدخول: المعي. والذي ضاعفه: ثلاث ساعات
دافئة في السيارة --- قاعدة \omterm{def:g6:exploring-our-environment:conditions}{الشروط} تغذي تضاعفًا كل نصف
ساعة تقريبًا. وأرخص سد: سلسلة البرد --- صندوق مبرّد (أو
أكله طازجًا)؛ وأرخص منه بعد ذلك: طهو دقيق مسبقًا.
\end{solution}

\begin{solution}{exo:g9:microbes-and-infection:10}
لأن الأيدي سعاة بين المقصورات: فأطباء سيميلفايس كانوا
نظافًا بأي معيار يومي، ومع ذلك حملوا قاتل الجناح من
الجثة إلى الأم. والغسلة التي تهم هي التي تكون \emph{عند
الحد} --- من الدجاج النيئ إلى السلطة، ومن المرحاض إلى
المائدة --- حيث كان الطريق سيصل لولاها.
\end{solution}

\begin{solution}{exo:g9:microbes-and-infection:11}
المضاد الحيوي يهاجم التضاعف --- فيكسر \omterm{def:g9:microbes-and-infection:sorted}{البكتيريا}
المنقسمة. وإن استعمل مهلهلًا صار كابح انتقاء يجري نصف
جريان: فتموت القابلة للتأثر، وترث الأقلية المقاومة
الجناح --- فيربّى الممرض الذي يلقاه المريض التالي.
\end{solution}

\begin{solution}{exo:g9:microbes-and-infection:12}
النص: قصاصة من أحرف \omterm{def:g9:heredity-and-traits:heredity}{وراثية}، بالأبجدية نفسها التي لكل
\omterm{def:g5:first-look-at-cells:cell}{خلية} (\cref{rem:g9:genes-and-dna:universal}). والمطبعة:
آلة القراءة والبناء في \omterm{def:g5:first-look-at-cells:cell}{خلية} حية، يحولها النص الداخل إلى
طبع نسخ من نفسه. وحافة الحياة: \omterm{def:g9:microbes-and-infection:sorted}{فالفيروس} وحده لا يفعل
شيئًا من قائمة \cref{prop:g1:living-or-not:signs} ---
لا تغذيًا، ولا نموًا، ولا انقسامًا؛ فإذا استعار جرى
محاكاة ساخرة للتكاثر. وصعوبة تدويته: إذ لا آلة فيروسية
تسمَّم --- بل المطبعة فقط، والمطبعة أنت.
\end{solution}

\begin{solution}{pb:g9:microbes-and-infection:1}
\textbf{1.} انتقال بالرذاذ --- سعالًا وكلامًا --- بمعونة
الأسطح المشتركة عند مضارب كرة الطاولة؛ والدخول عبر
المسالك الهوائية.

\textbf{2.} حساب التضاعف: \omterm{def:g5:first-look-at-cells:cell}{فخلايا} بطانة كل حالة المخطوفة
تطبع نسخًا تبذر الحالات التالية --- والتضاعف هو آلة
\omterm{def:g9:microbes-and-infection:sorted}{الفيروس} المستعارة، وهي تتراكم.

\textbf{3.} المرض. من ضرر \omterm{def:g5:first-look-at-cells:cell}{خلايا} البطانة المخطوفة
المخفقة --- وجزئيًا من معركة الدفاع نفسه، كما يفصّل
الفصل التالي.

\textbf{4.} لأن الاجتماعات توصل مدى رذاذ كل صف بالآخر
ساعة كاملة: فالانتقال يتدرج مع الازدحام، وأكثف خلطتين في
المدرسة كانتا الطريق بين المقصورات.

\textbf{5.} التلاميذ المرضى إلى بيوتهم: الانتقال مزال
عند المنبع. والنوافذ: جرعات الرذاذ مخففة --- انتقال.
وغسل الأيدي: طريق الأسطح مقطوع --- انتقال. ومسح
المعدات: طريق الأشياء المشتركة مقطوع --- انتقال. (وكلها
مصوَّبة إلى أرخص مرحلة نفسها.)

\textbf{6.} ``المضادات الحيوية تقتل \omterm{def:g9:microbes-and-infection:sorted}{البكتيريا}؛ وهذا
العامل \omterm{def:g9:microbes-and-infection:sorted}{فيروس} يجري على \omterm{def:g5:first-look-at-cells:cell}{خلايا} الأطفال أنفسهم --- فلا
يستطيع الدواء أن يمسه، ونثره مع ذلك لن يفعل إلا أن يجري
غربال انتقاء على \omterm{def:g9:microbes-and-infection:sorted}{بكتيريا} الجميع المقيمة، فيربّي المقاومة
ليوم تلزم فيه المضادات حقًا.''

\textbf{7.} العامل: \omterm{def:g9:microbes-and-infection:sorted}{بكتيريا} التسمم الغذائي. والانتقال:
الحلوى. والدخول: المعي. والتضاعف: ساعات بلا تبريد ---
دفء \omterm{def:g6:exploring-our-environment:conditions}{ورطوبة} وغذاء ووقت، كلها ممنوحة. والسد: الثلاجة ---
\omterm{def:g6:exploring-our-environment:conditions}{الشروط} مرفوضة؛ والفاشيتان تشتركان في مدرسة ولا شيء غير
ذلك.

\textbf{8.} خريطة الفاشية هي الانتقال وقد صار مرئيًا:
فمن-جلس-أين يرسم الطريق --- الصفوف، والرفاق،
والاجتماعات --- والطريق، متى رسم، بيّن أين يقطع.
فالعلاج العشوائي يحارب الحالات؛ أما الخريطة فتحارب
الانتشار.

\textbf{9.} الانتقال: رذاذ حالة الدليل، الصف والرفيقات
أولًا، ثم الاجتماعات إلى المدرسة. والدخول: المسالك
الهوائية، صفًا صفًا. والتضاعف: مضاعفات كل يومين خلال
الأسبوع الأول. والمرض والخمود: تبلغ الحميات ذروتها،
وتقطع الإجراءات الطرق، ويرق حوض القابلين --- فينحدر
المنحنى.

\textbf{10.} سدود الإجراءات جوّعت الطرق --- والحوض نفسه
كان أصغر من كشف الأسماء: فالتلاميذ الحاملون مناعة سلالة
الشتاء الماضي القريبة لم يمكن إشعالهم، والفاشية تموت
متى قلّ من يمكن أن يلتقطها: ذاكرة الفصل التالي، معروضة
سلفًا.

\textbf{11.} درس سيميلفايس: اقطع طرق السعاة --- العزل في
البيت، والأيدي المغسولة، والمضارب الممسوحة. ودرس
باستور: اعرف المسافر --- سمّ العامل، وارفض عليه \omterm{def:g6:exploring-our-environment:conditions}{الشروط}،
وصوّب العلاج إلى نوعه لا إلى العادة.

\textbf{12.} ``إيقاف الفاشيات أرخص ما يكون عند الانتقال
--- فالصابون والهواء والمسافة والصندوق المبرّد المغلق
تغلب أي دواء ينتظر التضاعف.''
\end{solution}
